\section*{Chapitre \ref{ch:hf:futures-market-making} --- Market making sur contrats à terme}

\begin{solution}{exo:hf:futures-market-making:1}
Priorité temporelle~: 100 et 150 aux deux premiers. Allocation au prorata~: $250\times$ pour chaque part de 1\,000, arrondie par défaut~: 25, 100, 75, 50. Allocation au premier ordre puis au prorata~: 100
au premier ordre, puis 150 répartis entre 900 lots~: 66, 50, 33 arrondis par défaut, et le lot restant au plus ancien, 67.
\end{solution}

\begin{solution}{exo:hf:futures-market-making:2}
Le carnet contient 2\,300 lots~: un agresseur de 1\,000 lots donne à chacun $1\,000\times230/2\,300=100$~; un agresseur de 7\,907 épuise le carnet et donne à chacun
sa pleine part de 230.
\end{solution}

\begin{solution}{exo:hf:futures-market-making:3}
Spread implicite à l'achat $9\,612-9\,599=13$, à la vente $9\,613-9\,598=15$. Acheter le spread à 12, vendre le contrat proche à 9\,612 et acheter le lointain à 9\,599~:
un tick, moins trois frais.
\end{solution}

\begin{solution}{exo:hf:futures-market-making:4}
La meilleure réponse de chaque teneur de marché aux tailles des autres est d'afficher davantage, quoi qu'ils fassent, mais lorsque tous le font, les parts restent inchangées et chaque teneur de marché supporte
davantage de risque sur les gros ordres~: utilité 19.5 à l'équilibre contre 22.3 si tous affichaient 100, avec dix teneurs de marché.
\end{solution}

\begin{solution}{exo:hf:futures-market-making:5}
Les teneurs de marché affichent bien plus qu'ils ne veulent négocier et annulent l'essentiel dès que le risque d'une grosse exécution augmente~; la taille affichée est un droit sur une part,
et non une intention de tout négocier.
\end{solution}

\begin{solution}{exo:hf:futures-market-making:6}
Il récompense l'amélioration du prix plutôt que l'affichage de taille~: le premier ordre à un nouveau prix est exécuté en premier, si bien que les teneurs de marché se disputent la fixation du prix avec
une taille modeste, ce qui modère la course à la taille aux prix existants.
\end{solution}

\begin{solution}{exo:hf:futures-market-making:7}
Avec la priorité temporelle, un teneur de marché afficherait 50 lots~; dix teneurs de marché en affichent 73, soit un facteur de 1.5 au lieu de 2.3~: une \omterm{def:hf:inventory-models:penalty}{pénalité d'inventaire} plus lourde freine la
spirale.
\end{solution}

\begin{solution}{exo:hf:futures-market-making:8}
L'essentiel des 20\,000 lots relève de la \omterm{def:hf:futures-market-making:over}{surcotation} (over-quoting), annulée dès qu'un gros ordre paraît probable~; la profondeur qui subsiste pour un ordre de 500 lots est
ce que les teneurs de marché veulent négocier, soit une fraction de ce qui est affiché.
\end{solution}

\begin{solution}{pb:hf:futures-market-making:1}
\begin{enumerate}
\item Voir le livre~1, chapitre~19, et l'encadré daté~: parts au prorata arrondies par défaut, lots en excédent servis dans l'ordre d'arrivée (premier entré, premier sorti), un ordre au premier rang servi d'abord lorsqu'il est prévu.
\item 100/150~; 25/100/75/50~; 100/67/50/33.
\item Voir la \cref{def:hf:futures-market-making:over}.
\item Les spreads valent presque toujours un tick, la profondeur dépasse de deux ordres de grandeur l'ordre au marché moyen, et les taux d'annulation dépassent 96~\%.
\item $f_i=s_i\min(X,S)/S$~; utilité $\mathrm{E}[hf_i-\gamma f_i^2/2]$.
\item 100 lots~: la taille à laquelle le gain espéré d'un lot marginal égale son coût espéré d'inventaire.
\item Meilleures réponses au total des autres, amorties, jusqu'à stabilisation.
\item Avec de nombreux teneurs de marché, la part de chacun est faible et une taille plus grande achète une part à bon marché~: le risque d'épuiser le carnet ne rattrape jamais
avant le plafond.
\item 1.1, 2.3, 6.2 et 50 (au plafond).
\item Plus faible pour chaque $n$~: 19.5 contre 22.3 avec dix teneurs de marché, 12.3 contre 14.3 avec vingt.
\item Exécute chacun des dix teneurs de marché à hauteur de 230 lots, soit 2.3 fois la taille souhaitée, tous du même côté.
\item Voir la \cref{def:hf:futures-market-making:implied}~; 0 pour un décalage d'un événement, 0.9, 5.9, 23.7 et 49.7 pour mille pour 5, 10, 20 et 30.
\item Avec dix teneurs de marché, chacun affiche 2.3 fois sa taille souhaitée~; avec vingt, ils atteignent le maximum autorisé par la bourse~; un ordre au 99.9\up{e}~centile laisse chaque teneur de marché
2.3 fois son inventaire souhaité d'un seul coup.
\item À la taille d'équilibre, avec une annulation rapide lorsque les gros ordres deviennent probables, et une limite sur l'exécution qu'il peut supporter.
\item Allocation au premier ordre ou pondérée par le temps, une durée minimale de présence dans le carnet, ou des frais sur la taille affichée mais non négociée.
\item Cohérence~: le prix de chaque carnet provient d'une même courbe, si bien que le teneur de marché n'est jamais victime d'arbitrage entre ses propres cotations.
\item Lorsque la relation entre les deux matières premières change (panne d'une raffinerie, nouvelle règle de livraison).
\item \omterm{def:hf:futures-market-making:implied}{Arbitrage de prix implicite}.
\item Modéliser les tailles des autres à partir de la profondeur totale du carnet et simuler l'allocation avec \texttt{firm.match}.
\item Une part de chaque ordre, et avec elle le risque de chaque gros ordre.
\end{enumerate}
\end{solution}

\begin{solution}{iq:hf:futures-market-making:1}
Le tick est grand par rapport à la volatilité du contrat sur l'horizon d'une transaction, et de nombreux teneurs de marché se font concurrence~: la fourchette ne peut être plus étroite qu'un tick,
et la concurrence l'empêche d'être plus large.
\iqlookfor{tick large, concurrence.}
\end{solution}

\begin{solution}{iq:hf:futures-market-making:2}
Avec $f_i=Xs_i/S$~: $\mathrm{E}[f_i]=\mu s_i/S$, $\mathrm{E}[f_i^2]=m_2s_i^2/S^2$~; maximiser $h\mu p-\gamma m_2p^2/2$ sur la part $p=s_i/S$~:
$p^\ast=h\mu/(\gamma m_2)$, d'où $s_i=p^\ast S_{-i}/(1-p^\ast)$ si $p^\ast<1$. Dans un équilibre symétrique $p=1/n$, donc si $1/n<p^\ast$ chaque teneur de marché
veut davantage et les tailles croissent sans limite~: le terme d'épuisement du carnet est ce qui l'arrête.
\iqlookfor{parts, et pourquoi le cas non borné survient.}
\end{solution}

\begin{solution}{iq:hf:futures-market-making:3}
On se fait exécuter à hauteur de 2\,000 fois sa part du carnet~; chaque teneur de marché est exécuté d'un coup et le prix franchit probablement le niveau. Annuler tous les ordres au repos
de ce côté, couvrir sur le contrat lié le plus liquide, et réévaluer avant de coter à nouveau.
\iqlookfor{exécutions corrélées et réduction immédiate du risque.}
\end{solution}

\begin{solution}{iq:hf:futures-market-making:4}
Parts par partie entière de la quantité multipliée par la taille, rapportée au total~; écarter les parts inférieures au minimum~; les restes par ordre temporel. Cas limites~: agresseur plus grand que le carnet, toutes les parts arrondies à zéro,
ordres entièrement exécutés avant la passe sur les restes, égalités temporelles.
\iqlookfor{arrondis, minimums, épuisement.}
\end{solution}

\begin{solution}{iq:hf:futures-market-making:5}
Il y a infraction si une exécution complète dépasserait la limite de position~: les contrôles de risque pré-négociation doivent compter la taille affichée comme position potentielle~; la firme a besoin d'une
limite sur la taille affichée ou d'une garantie d'annulation.
\iqlookfor{position potentielle.}
\end{solution}

\begin{solution}{iq:hf:futures-market-making:6}
Comparer la profondeur à la distribution des tailles des agresseurs et aux annulations précédant les grosses transactions~; modéliser les tailles souhaitées par les teneurs de marché à partir de leur
comportement dans des carnets à priorité temporelle de contrats similaires.
\iqlookfor{profondeur par rapport à la taille exécutée et comportement d'annulation.}
\end{solution}
